\documentclass[11pt,a4paper]{article}

\usepackage[T1]{fontenc}
\usepackage[utf8]{inputenc}
\usepackage{lmodern}
\usepackage{microtype}
\usepackage[a4paper,margin=30mm]{geometry}
\usepackage{amsmath,amssymb,amsthm,mathtools}
\usepackage{mathrsfs}
\usepackage{booktabs,array,tabularx}
\usepackage{enumitem}
\usepackage[hidelinks]{hyperref}
\hypersetup{
  pdftitle={An Exact Divisor-Theoretic Classification of the Integer Solutions of z squared plus y squared z plus x cubed minus 2 equals 0},
  pdfauthor={Jamal Agbanwa},
  pdfsubject={A complete divisor-theoretic classification of the integer solutions},
  pdfkeywords={Diophantine equation, integer solutions, divisor parametrization}
}
\usepackage{fancyhdr}

\allowdisplaybreaks
\setlength{\emergencystretch}{2em}

\pagestyle{fancy}
\fancyhf{}
\fancyhead[L]{Integer solutions of $z^2+y^2z+x^3-2=0$}
\fancyhead[R]{J. Agbanwa}
\fancyfoot[C]{\thepage}
\setlength{\headheight}{14pt}

\newtheorem{theorem}{Theorem}[section]
\newtheorem{proposition}[theorem]{Proposition}
\newtheorem{lemma}[theorem]{Lemma}
\newtheorem{corollary}[theorem]{Corollary}
\theoremstyle{definition}
\newtheorem{definition}[theorem]{Definition}
\theoremstyle{remark}
\newtheorem{remark}[theorem]{Remark}

\newcommand{\Z}{\mathbb Z}
\newcommand{\N}{\mathbb N}
\newcommand{\Sols}{\mathscr S}
\newcommand{\Yset}{\mathcal Y}
\newcommand{\eqE}{\textup{(E)}}

\title{\textbf{An Exact Divisor-Theoretic Classification of the Integer Solutions of}
\\[3mm]
$z^2+y^2z+x^3-2=0$}
\author{}
\date{}

\begin{document}

\maketitle

\begin{abstract}
We give a complete elementary classification of the triples
$(x,y,z)\in\Z^3$ satisfying
\[
 z^2+y^2z+x^3-2=0.
\]
The central observation is not merely a construction but an explicit
bijection: a solution corresponds uniquely to an ordered pair of integer
factors $a,b$ satisfying
$ab=x^3-2$ and $a+b=y^2$, through
$a=-z$ and $b=z+y^2$.  Resolving the signs of these two factors yields a
canonical description in two exhaustive cases, $x\geq 2$ and $x\leq 1$.
We prove both directions of the classification, its uniqueness statement,
the treatment of all diagonal and zero-square cases, a compact ordered-divisor
formulation, a root-swapping involution, and a terminating divisor algorithm
which enumerates all solutions for each fixed value of $x$.
\end{abstract}

\medskip
\noindent\textbf{Keywords.}
Diophantine equation; integer solutions; divisor parametrization; quadratic
factorization; integral points.

\section{Introduction and notation}

We study the Diophantine equation
\begin{equation*}
 z^2+y^2z+x^3-2=0,
 \qquad (x,y,z)\in\Z^3.
 \tag{E}
\end{equation*}
The purpose of this note is to describe its entire integer solution set by an
exact divisor criterion.  The word ``exact'' is important: every triple
produced by the criterion is verified to satisfy \eqE, and every
integer solution of \eqE is proved to arise from the criterion with
uniquely determined factor data, up to the explicitly stated coincidences.

We use
\[
 \N=\{1,2,3,\ldots\},
 \qquad
 \Z_{\geq 0}=\{0,1,2,\ldots\}.
\]
For $s\in\Z_{\geq0}$, define
\begin{equation}
 \Yset(s)=
 \begin{cases}
  \{0\},&s=0,\\
  \{-s,s\},&s>0.
 \end{cases}
 \label{eq:Yset}
\end{equation}
Thus $\Yset(s)$ is exactly the set of integers whose absolute value is $s$;
the definition avoids counting $+0$ and $-0$ as different values.

Let
\[
 \Sols=\{(x,y,z)\in\Z^3:z^2+y^2z+x^3-2=0\}
\]
and, for a fixed $x\in\Z$, let
\[
 \Sols_x=\{(x,y,z)\in\Sols\}.
\]

The factorization driving the entire argument is
\begin{equation}
 (-z)(z+y^2)=x^3-2.
 \label{eq:factorization}
\end{equation}
The main work is to formulate this identity as a bijection and then resolve
all possible signs without losing or duplicating solutions.

\section{The exact factor-pair correspondence}

We begin by ruling out a degenerate value which could otherwise make one of
the factors in \eqref{eq:factorization} vanish.

\begin{lemma}\label{lem:no-cube-two}
For every integer $x$, one has $x^3\neq2$.
Consequently, if $(x,y,z)\in\Sols$, then
\[
 z\neq0
 \qquad\text{and}\qquad
 z+y^2\neq0.
\]
\end{lemma}

\begin{proof}
If $x\leq1$, then $x^3\leq1<2$.  If $x\geq2$, then
$x^3\geq8>2$.  Since every integer lies in exactly one of these two cases,
$x^3=2$ has no integer solution.

Now suppose $(x,y,z)\in\Sols$.  Rearranging \eqE gives
\[
 z(z+y^2)=2-x^3.
\]
By the first part, the right-hand side is nonzero.  A product of two integers
is nonzero only when both factors are nonzero, so $z\neq0$ and
$z+y^2\neq0$.
\end{proof}

The next proposition is the precise logical core of the classification.

\begin{proposition}[Factor-pair bijection]\label{prop:bijection}
Let
\[
 \mathscr T=
 \bigl\{(x,y,a,b)\in\Z^4:
       ab=x^3-2\ \text{and}\ a+b=y^2\bigr\}.
\]
Define
\begin{align*}
 \Phi:\Sols&\longrightarrow\mathscr T,
 &\Phi(x,y,z)&=(x,y,-z,z+y^2),\\
 \Psi:\mathscr T&\longrightarrow\Sols,
 &\Psi(x,y,a,b)&=(x,y,-a).
\end{align*}
Then $\Phi$ and $\Psi$ are well-defined mutually inverse bijections.
In particular, for every solution $(x,y,z)$, the associated ordered factor
pair is uniquely
\begin{equation}
 a=-z,
 \qquad
 b=z+y^2.
 \label{eq:unique-ab}
\end{equation}
\end{proposition}

\begin{proof}
First let $(x,y,z)\in\Sols$ and put
$a=-z$ and $b=z+y^2$.  Then
\[
 a+b=-z+z+y^2=y^2.
\]
Moreover, using \eqE,
\[
 ab=(-z)(z+y^2)
    =-z^2-y^2z
    =x^3-2.
\]
Thus $(x,y,a,b)\in\mathscr T$, so $\Phi$ is well-defined.

Conversely, let $(x,y,a,b)\in\mathscr T$ and put $z=-a$.
Since $a+b=y^2$ and $ab=x^3-2$, we obtain
\begin{align*}
 z^2+y^2z+x^3-2
 &=a^2-a(a+b)+ab\\
 &=a^2-a^2-ab+ab\\
 &=0.
\end{align*}
Hence $(x,y,-a)\in\Sols$, so $\Psi$ is well-defined.

It remains to verify that the two maps are inverse to one another.  For
$(x,y,z)\in\Sols$,
\[
 (\Psi\circ\Phi)(x,y,z)
 =\Psi(x,y,-z,z+y^2)
 =(x,y,z).
\]
For $(x,y,a,b)\in\mathscr T$, the relation $a+b=y^2$ gives
\begin{align*}
 (\Phi\circ\Psi)(x,y,a,b)
 &=\Phi(x,y,-a)\\
 &=(x,y,a,-a+y^2)\\
 &=(x,y,a,b).
\end{align*}
Therefore $\Phi$ and $\Psi$ are mutually inverse bijections.  Formula
\eqref{eq:unique-ab} follows directly from the definition of $\Phi$ and also
shows uniqueness of the ordered factor pair.
\end{proof}

\section{Sign analysis and direct verification}

The sign restrictions needed below are elementary, but we state and prove
them explicitly so that no case is hidden.

\begin{lemma}[Positive product]\label{lem:positive-product}
Let $a,b\in\Z$.  If $ab>0$ and $a+b\geq0$, then $a>0$ and $b>0$.
\end{lemma}

\begin{proof}
The inequality $ab>0$ implies that $a$ and $b$ have the same sign and are
both nonzero.  If they were both negative, then $a+b<0$, contradicting the
hypothesis.  Hence both are positive.
\end{proof}

\begin{lemma}[Negative product]\label{lem:negative-product}
Let $a,b\in\Z$.  If $ab<0$ and $a+b\geq0$, then there exist unique integers
$d,e\in\N$ with $d\leq e$ such that
\[
 \{a,b\}=\{-d,e\}.
\]
For these integers,
\[
 |ab|=de,
 \qquad
 a+b=e-d.
\]
\end{lemma}

\begin{proof}
Since $ab<0$, exactly one of $a,b$ is positive and the other is negative.
Let $e$ be the positive one and write the negative one as $-d$, where
$d,e\in\N$.  Then
\[
 a+b=e-d.
\]
The assumption $a+b\geq0$ implies $e-d\geq0$, hence $d\leq e$.
Furthermore,
\[
 |ab|=|-de|=de.
\]
The positive member of the set $\{a,b\}$ is uniquely determined, as is the
absolute value of its negative member; therefore $d$ and $e$ are unique.
\end{proof}

We shall use the following direct substitution in the soundness part of the
main theorem.

\begin{lemma}[Direct verification of the two constructions]
\label{lem:direct-verification}
Let $x\in\Z$ and let $d,e\in\N$.

\begin{enumerate}[label=\textup{(\alph*)}]
\item If
\[
 de=x^3-2,
 \qquad
 d+e=s^2
\]
for some $s\in\Z_{\geq0}$, then, for every $y\in\Yset(s)$, both
\[
 (x,y,-d)
 \qquad\text{and}\qquad
 (x,y,-e)
\]
satisfy \eqE.

\item If
\[
 de=2-x^3,
 \qquad
 e-d=s^2
\]
for some $s\in\Z_{\geq0}$, then, for every $y\in\Yset(s)$, both
\[
 (x,y,d)
 \qquad\text{and}\qquad
 (x,y,-e)
\]
satisfy \eqE.
\end{enumerate}
\end{lemma}

\begin{proof}
In each case $y^2=s^2$ by the definition of $\Yset(s)$.

For part (a), if $z=-d$, then
\begin{align*}
 z^2+y^2z+x^3-2
 &=d^2-d(d+e)+de\\
 &=d^2-d^2-de+de\\
 &=0.
\end{align*}
If $z=-e$, then similarly
\begin{align*}
 z^2+y^2z+x^3-2
 &=e^2-e(d+e)+de\\
 &=e^2-ed-e^2+de\\
 &=0.
\end{align*}

For part (b), the identity $de=2-x^3$ is equivalent to
$x^3-2=-de$.  If $z=d$, then
\begin{align*}
 z^2+y^2z+x^3-2
 &=d^2+d(e-d)-de\\
 &=d^2+de-d^2-de\\
 &=0.
\end{align*}
If $z=-e$, then
\begin{align*}
 z^2+y^2z+x^3-2
 &=e^2-e(e-d)-de\\
 &=e^2-e^2+ed-de\\
 &=0.
\end{align*}
This proves both assertions.
\end{proof}

\section{The complete canonical classification}

We now give the promised description of all integer solutions.

\begin{theorem}[Complete classification]\label{thm:classification}
The integer solutions of \eqE are exactly the triples described in
one of the following two mutually exclusive cases.

\medskip
\noindent\textup{\textbf{I. The case $x\geq2$.}}
Choose integers $d,e,s$ satisfying
\begin{equation}
 d,e\in\N,
 \qquad
 s\in\Z_{\geq0},
 \qquad
 d\leq e,
 \qquad
 de=x^3-2,
 \qquad
 d+e=s^2.
 \label{eq:case-positive-data}
\end{equation}
Then $s>0$, and for every $y\in\Yset(s)$ the triples
\begin{equation}
 (x,y,-d)
 \qquad\text{and}\qquad
 (x,y,-e)
 \label{eq:case-positive-solutions}
\end{equation}
are solutions.  Conversely, every solution with $x\geq2$ is of this form.
For a given solution, the canonical data $(d,e,s)$ are unique.  If $d=e$,
the two displayed values of $z$ in \eqref{eq:case-positive-solutions}
coincide and are to be retained only once.

\medskip
\noindent\textup{\textbf{II. The case $x\leq1$.}}
Choose integers $d,e,s$ satisfying
\begin{equation}
 d,e\in\N,
 \qquad
 s\in\Z_{\geq0},
 \qquad
 d\leq e,
 \qquad
 de=2-x^3,
 \qquad
 e-d=s^2.
 \label{eq:case-negative-data}
\end{equation}
Then for every $y\in\Yset(s)$ the triples
\begin{equation}
 (x,y,d)
 \qquad\text{and}\qquad
 (x,y,-e)
 \label{eq:case-negative-solutions}
\end{equation}
are solutions.  Conversely, every solution with $x\leq1$ is of this form.
For a given solution, the canonical data $(d,e,s)$ are unique.  When $d=e$,
one has $s=0$, so the only value of $y$ is $0$; the two values of $z$ are
$d$ and $-d$, which are distinct because $d>0$.
\end{theorem}

\begin{proof}
We prove soundness, completeness, and uniqueness in each case.

\smallskip
\noindent\emph{Case I: $x\geq2$.}
Here $x^3-2\geq6>0$.

Suppose first that data $d,e,s$ satisfy
\eqref{eq:case-positive-data}.  The asserted triples are solutions by
Lemma~\ref{lem:direct-verification}(a).  Since $d,e\in\N$, one has
$d+e\geq2$, and therefore $s^2=d+e>0$; hence $s>0$.
This proves soundness.

Conversely, let $(x,y,z)\in\Sols$ with $x\geq2$.  By
Proposition~\ref{prop:bijection}, define
\[
 a=-z,
 \qquad
 b=z+y^2.
\]
Then
\[
 ab=x^3-2>0,
 \qquad
 a+b=y^2\geq0.
\]
Lemma~\ref{lem:positive-product} gives $a,b>0$.  Put
\[
 d=\min\{a,b\},
 \qquad
 e=\max\{a,b\},
 \qquad
 s=|y|.
\]
Then $d,e\in\N$, $d\leq e$, and
\[
 de=ab=x^3-2,
 \qquad
 d+e=a+b=y^2=s^2.
\]
Moreover, $z=-a$, so $z=-d$ or $z=-e$.  Thus the given solution appears in
\eqref{eq:case-positive-solutions}.  This proves completeness.

For uniqueness, Proposition~\ref{prop:bijection} shows that the ordered pair
$(a,b)=(-z,z+y^2)$ is uniquely determined by the solution.  Hence its
canonically ordered version $(d,e)=(\min\{a,b\},\max\{a,b\})$ is unique, and
$s=|y|$ is unique.  If $d=e$, then $-d=-e$, so the two displayed choices of
$z$ are the same; if $d<e$, they are distinct.

\smallskip
\noindent\emph{Case II: $x\leq1$.}
Here $x^3-2\leq-1<0$, or equivalently $2-x^3>0$.

Suppose first that data $d,e,s$ satisfy
\eqref{eq:case-negative-data}.  The asserted triples are solutions by
Lemma~\ref{lem:direct-verification}(b).  This proves soundness.

Conversely, let $(x,y,z)\in\Sols$ with $x\leq1$, and again put
\[
 a=-z,
 \qquad
 b=z+y^2.
\]
Proposition~\ref{prop:bijection} gives
\[
 ab=x^3-2<0,
 \qquad
 a+b=y^2\geq0.
\]
By Lemma~\ref{lem:negative-product}, there exist unique $d,e\in\N$ with
$d\leq e$ and
\[
 \{a,b\}=\{-d,e\}.
\]
It follows that
\[
 de=|ab|=2-x^3,
 \qquad
 e-d=a+b=y^2.
\]
With $s=|y|$, we therefore have $e-d=s^2$.  Since $z=-a$, there are precisely
two possibilities:
\[
 a=-d\ \Longrightarrow\ z=d,
 \qquad
 a=e\ \Longrightarrow\ z=-e.
\]
Thus the given solution appears in \eqref{eq:case-negative-solutions}, proving
completeness.

Uniqueness of $d,e$ follows from Lemma~\ref{lem:negative-product}, or directly
from the unique ordered pair $(a,b)$ supplied by
Proposition~\ref{prop:bijection}; again $s=|y|$ is unique.  If $d=e$, then
$s^2=e-d=0$, so $s=0$ and $y=0$.  The corresponding values $z=d$ and
$z=-e=-d$ are distinct because $d\in\N$.

Finally, every integer $x$ satisfies exactly one of $x\geq2$ and $x\leq1$.
Thus the two cases together are exhaustive and mutually exclusive.
\end{proof}

\begin{remark}[Logical content of the theorem]\label{rem:logical-content}
Theorem~\ref{thm:classification} contains all three ingredients required for
a genuine classification:

\begin{enumerate}[label=\textup{(\roman*)},leftmargin=2.5em]
\item every listed triple satisfies the original equation;
\item every integer solution is recovered by the construction;
\item the canonical factor data attached to a solution are unique, with all
possible presentational coincidences explicitly identified.
\end{enumerate}
No unproved existence assumption is used in any direction.
\end{remark}

\section{Equivalent compact formulations and symmetry}

The canonical form of Theorem~\ref{thm:classification} is convenient for sign
analysis.  The following ordered-divisor form is more compact.

\begin{corollary}[Ordered-divisor parametrization]\label{cor:ordered-divisor}
A triple $(x,y,z)\in\Z^3$ satisfies \eqE if and only if there exist
$a\in\Z\setminus\{0\}$ and $s\in\Z_{\geq0}$ such that
\begin{equation}
 a\mid x^3-2,
 \qquad
 a+\frac{x^3-2}{a}=s^2,
 \qquad
 y\in\Yset(s),
 \qquad
 z=-a.
 \label{eq:ordered-divisor-form}
\end{equation}
For a given solution, the parameter $a$ is uniquely $a=-z$, and
$s=|y|$.
\end{corollary}

\begin{proof}
Suppose $(x,y,z)$ is a solution.  By Lemma~\ref{lem:no-cube-two}, $z\neq0$,
so $a=-z$ is nonzero.  Proposition~\ref{prop:bijection} gives
\[
 a\left(z+y^2\right)=x^3-2,
\]
so $a\mid x^3-2$, and
\[
 z+y^2=\frac{x^3-2}{a}.
\]
As $z=-a$, this becomes
\[
 a+\frac{x^3-2}{a}=y^2.
\]
Taking $s=|y|$ gives \eqref{eq:ordered-divisor-form}.

Conversely, suppose \eqref{eq:ordered-divisor-form} holds.  Put
$b=(x^3-2)/a$.  Then $ab=x^3-2$ and $a+b=s^2=y^2$.  The inverse map
$\Psi$ in Proposition~\ref{prop:bijection} therefore shows that
$(x,y,-a)$ is a solution.  Uniqueness of $a=-z$ and $s=|y|$ is immediate.
\end{proof}

For fixed $x$ and $y$, the equation is quadratic in $z$.  Its two roots are
interchanged by a simple involution.

\begin{proposition}[Root-swapping involution]\label{prop:involution}
Define
\[
 \iota:\Sols\longrightarrow\Sols,
 \qquad
 \iota(x,y,z)=(x,y,-y^2-z).
\]
Then $\iota$ is an involution: $\iota^2$ is the identity on $\Sols$.
Under the bijection of Proposition~\ref{prop:bijection}, $\iota$ exchanges the
ordered factors $a$ and $b$.

A solution is fixed by $\iota$ if and only if
\[
 2z+y^2=0,
\]
equivalently if and only if its associated factors satisfy $a=b$.
Such a fixed point can occur only in the case $x\geq2$.
\end{proposition}

\begin{proof}
Let $(x,y,z)\in\Sols$ and put $z'=-y^2-z$.  Then
\begin{align*}
 (z')^2+y^2z'
 &=(-y^2-z)^2+y^2(-y^2-z)\\
 &=(y^2+z)^2-y^2(y^2+z)\\
 &=z(y^2+z)\\
 &=z^2+y^2z.
\end{align*}
Therefore
\[
 (z')^2+y^2z'+x^3-2
 =z^2+y^2z+x^3-2
 =0,
\]
so $\iota$ maps $\Sols$ to itself.  Applying the map twice gives
\[
 -y^2-(-y^2-z)=z,
\]
so $\iota^2$ is the identity.

For the factor correspondence, if
$a=-z$ and $b=z+y^2$, then for $z'=-y^2-z$ one has
\[
 -z'=y^2+z=b,
 \qquad
 z'+y^2=-z=a.
\]
Thus $a$ and $b$ are exchanged.

The fixed-point condition $z' = z$ is exactly
$-y^2-z=z$, or $2z+y^2=0$.  Since $\iota$ exchanges $a$ and $b$, this is also
equivalent to $a=b$.  But $a=b$ implies $ab=a^2>0$ because
$a\neq0$ by Lemma~\ref{lem:no-cube-two}.  Hence $x^3-2=ab>0$, which forces
$x\geq2$.
\end{proof}

\section{Effective enumeration for fixed \texorpdfstring{$x$}{x}}

The classification immediately yields a finite, terminating procedure for
each fixed integer $x$.

\begin{proposition}[Complete divisor algorithm]\label{prop:algorithm}
For a fixed $x\in\Z$, all elements of $\Sols_x$ are obtained by the following
finite procedure.

\begin{enumerate}[leftmargin=2.2em]
\item If $x\geq2$, set $N=x^3-2$.  Enumerate every positive divisor $d$ of
$N$ satisfying $d\leq N/d$, and put $e=N/d$.  Retain the pair $(d,e)$ exactly
when $d+e$ is a square, say $d+e=s^2$.  For each retained pair output
\[
 (x,y,-d),\quad (x,y,-e)
 \qquad (y\in\Yset(s)),
\]
with repeated triples removed.

\item If $x\leq1$, set $N=2-x^3$.  Enumerate every positive divisor $d$ of
$N$ satisfying $d\leq N/d$, and put $e=N/d$.  Retain the pair $(d,e)$ exactly
when $e-d$ is a square, say $e-d=s^2$.  For each retained pair output
\[
 (x,y,d),\quad (x,y,-e)
 \qquad (y\in\Yset(s)).
\]
\end{enumerate}
The procedure terminates, outputs only solutions, and outputs every solution
with the prescribed value of $x$.
\end{proposition}

\begin{proof}
In either case, $N$ is a positive integer, so it has only finitely many
positive divisors.  The enumeration therefore terminates.  Every output is a
solution by Lemma~\ref{lem:direct-verification}.  Conversely, the completeness
part of Theorem~\ref{thm:classification} says that every solution in
$\Sols_x$ has exactly one of the retained canonical factor pairs.  Hence no
solution is omitted.
\end{proof}

For completeness, we record an exact counting formula in terms of admissible
factor pairs and a simple universal bound for each fiber.

\begin{corollary}[Finiteness and a divisor bound]\label{cor:counting}
For every $x\in\Z$, the set $\Sols_x$ is finite.  Define the set of admissible
canonical factor pairs by
\[
 \mathcal Q_x=
 \begin{cases}
 \{(d,e)\in\N^2:d\leq e,\ de=x^3-2,\ d+e\text{ is a square}\},
   &x\geq2,\\[1mm]
 \{(d,e)\in\N^2:d\leq e,\ de=2-x^3,\ e-d\text{ is a square}\},
   &x\leq1.
 \end{cases}
\]
Then
\begin{equation}
 \#\Sols_x
 =4\,\#\{(d,e)\in\mathcal Q_x:d<e\}
  +2\,\#\{(d,d)\in\mathcal Q_x\}.
 \label{eq:exact-count}
\end{equation}
If $\tau(n)$ denotes the number of positive divisors of $n$, then
\begin{equation}
 \#\Sols_x\leq 2\tau\bigl(|x^3-2|\bigr).
 \label{eq:divisor-bound}
\end{equation}
\end{corollary}

\begin{proof}
Finiteness follows at once from Proposition~\ref{prop:algorithm}.

Let $(d,e)\in\mathcal Q_x$.  If $d<e$, then the relevant square is positive:
for $x\geq2$ it is $d+e>0$, and for $x\leq1$ it is $e-d>0$.  Hence there are
two values of $y$.  There are also two distinct values of $z$, because
$-d\neq-e$ in the first case and $d\neq-e$ in the second.  Thus the pair
contributes exactly four solutions.

If $d=e$, then the pair contributes exactly two solutions.  Indeed, for
$x\geq2$ the two displayed values of $z$ coincide, while $d+e=2d>0$ gives
two values of $y$ whenever the pair is admissible.  For $x\leq1$, one has
$e-d=0$, so $y=0$ is the only value, while the two values of $z$ are $d$ and
$-d$.  The uniqueness assertion in Theorem~\ref{thm:classification} ensures
that contributions from distinct canonical pairs do not overlap.  This
proves \eqref{eq:exact-count}.

Put $n=|x^3-2|$.  The positive factor pairs $d\leq e$ with $de=n$ are the
orbits of the positive divisors under $d\mapsto n/d$.  If $n$ is not a
square, then $\tau(n)=2r$ and there are $r$ pairs, all with $d<e$; even if all
were admissible, \eqref{eq:exact-count} would give at most
$4r=2\tau(n)$ solutions.  If $n$ is a square, then $\tau(n)=2r+1$; there are
$r$ pairs with $d<e$ and one diagonal pair.  The maximum possible
contribution is then
\[
 4r+2=2(2r+1)=2\tau(n).
\]
This proves \eqref{eq:divisor-bound}.
\end{proof}

\section{Illustrative examples}

The following table displays selected admissible factor pairs and the
solutions contributed by each pair.  Each row is justified immediately by
Theorem~\ref{thm:classification}; no separate search in $z$ is needed.

\begin{center}
\renewcommand{\arraystretch}{1.25}
\begin{tabularx}{\textwidth}{@{}c>{\raggedright\arraybackslash}X>{\raggedright\arraybackslash}X>{\raggedright\arraybackslash}X@{}}
\toprule
$x$ & Factor identity & Square identity & Solutions contributed \\
\midrule
$1$
& $2-1^3=1\cdot1$
& $1-1=0^2$
& $(1,0,1)$ and $(1,0,-1)$ \\

$0$
& $2-0^3=1\cdot2$
& $2-1=1^2$
& $(0,\pm1,1)$ and $(0,\pm1,-2)$ \\

$-2$
& $2-(-2)^3=1\cdot10$
& $10-1=3^2$
& $(-2,\pm3,1)$ and $(-2,\pm3,-10)$ \\

$-7$
& $2-(-7)^3=5\cdot69$
& $69-5=8^2$
& $(-7,\pm8,5)$ and $(-7,\pm8,-69)$ \\

$8$
& $8^3-2=15\cdot34$
& $15+34=7^2$
& $(8,\pm7,-15)$ and $(8,\pm7,-34)$ \\

$20$
& $20^3-2=31\cdot258$
& $31+258=17^2$
& $(20,\pm17,-31)$ and $(20,\pm17,-258)$ \\
\bottomrule
\end{tabularx}
\end{center}

A single value of $x$ can admit more than one canonical factor pair.  For
example,
\[
 26^3-2=17574=58\cdot303=87\cdot202,
\]
and
\[
 58+303=361=19^2,
 \qquad
 87+202=289=17^2.
\]
Consequently these two factor pairs yield the eight solutions
\begin{align*}
 &(26,\pm19,-58),\qquad (26,\pm19,-303),\\
 &(26,\pm17,-87),\qquad (26,\pm17,-202).
\end{align*}
The uniqueness clause of Theorem~\ref{thm:classification} shows that the two
factor pairs account for disjoint sets of triples.

\section{Conclusion}

The substitution
\[
 a=-z,
 \qquad
 b=z+y^2
\]
turns \eqE into the simultaneous relations
\[
 ab=x^3-2,
 \qquad
 a+b=y^2,
\]
and Proposition~\ref{prop:bijection} proves that this transformation is a
bijection, not merely a necessary condition.  The sign of $x^3-2$ then forces
exactly the two canonical alternatives in
Theorem~\ref{thm:classification}.  Therefore the theorem gives an exhaustive,
sound, and unique divisor-theoretic description of every integer solution.
For each fixed $x$, the classification is effective and finite, as made
explicit by Proposition~\ref{prop:algorithm}.

The result should be understood precisely in this sense: it completely
classifies the triples through an exact factor-pair criterion and a finite
algorithm on every fiber $\Sols_x$.  It does not require, and does not claim,
a separate closed-form description of the set of all $x$ for which
$\Sols_x$ is nonempty.

\end{document}
